% Source-derived chapter 04: Coherent sheaves on $\sY$
% Original source: tates-thesis-de-rham-setting
\chapter{Coherent sheaves on $\sY$}
\label{tates-thesis-de-rham-setting-chapter-04}
\source{tates-thesis-de-rham-setting}

\begin{remark}[Source section]
\label{tates-thesis-de-rham-setting-chapter-04-source}
\source{tates-thesis-de-rham-setting}
\source{tates-thesis-de-rham-setting}
This chapter reproduces the corresponding section of the retrieved
LaTeX source, including its definitions, statements, examples, and
proofs. It has not yet been translated into Lean.
\notready
\end{remark}

% The source section title was: Coherent sheaves on $\sY$
\label{s:indcoh}
\source{tates-thesis-de-rham-setting}

\subsection{Overview}

In this section, we define $\IndCoh^*$ for $\sY$ and its relatives.
This section follows \cite{methods} \S 6, though we keep our exposition largely
independent of \emph{loc. cit}. 

We strive to keep our exposition as explicit as possible, so assume
whatever simplifying assumptions we like. We refer to \emph{loc. cit}. for
a more thorough development of the subject.

\subsection{Affine case}

\subsubsection{}

Suppose $S$ is an eventually coconnective (e.g. classical) affine scheme.

We define the (non-cocomplete, DG) subcategory $\Coh(S) \subset \QCoh(S)$
to consist of objects $\sF \in \QCoh(S)^+$ such that for each $n$, the object
$\tau^{\geq -n}(\sF) \in \QCoh(S)^{\geq -n}$ is compact. 
In other words, $\sF$ should be eventually coconnective and 
$\ul{\Hom}_{\QCoh(S)}(\sF,-)$ should commute with
filtered colimits in $\QCoh(S)^{\geq -n}$ for each $n$.
Clearly $\Perf(S) \subset \Coh(S)$.

\begin{notation}
\source{tates-thesis-de-rham-setting}
\notready

For $S = \Spec(A)$, we also sometimes write $\Coh(A) \subset A\mod$ for 
the subcategory $\Coh(S) \subset \QCoh(S) \simeq A\mod$.

\end{notation}

We define $\IndCoh^*(S)$ as $\Ind(\Coh(S))$. 
Observe that we have a
natural functor $\Psi:\IndCoh^*(S) \to \QCoh(S)$.

Define $\IndCoh^*(S)^{\leq 0}$ as the subcategory generated
under colimits by coherent objects that are connective (with respect to the
$t$-structure on $\QCoh(S)$). There is a corresponding 
$t$-structure on $\IndCoh^*(S)$ with this as its connective subcategory.

\begin{lem}[\cite{methods} Lemma 6.4.1]
\source{tates-thesis-de-rham-setting}
\notready\label{l:psi-affine}
\source{tates-thesis-de-rham-setting}

The above functor $\Psi$ is $t$-exact and induces an equivalence 
$\Psi:\IndCoh^*(S)^+ \isom \QCoh(S)^+$.

\end{lem}

\subsubsection{Pullbacks}

Let $f:S \to T$ be an eventually coconnective morphism of affine, eventually
coconnective schemes.


\subsubsection{Pushforwards}

Now suppose $f:S \to T$ is a morphism of affine, eventually coconnective schemes.
In this case, there is a unique continuous DG functor:
%
\[
f_*^{\IndCoh}:\IndCoh^*(S) \to \IndCoh^*(T)
\]

\noindent fitting into a commutative diagram:
%
\[
\xymatrix{
\IndCoh^*(S) \ar[rr]^{f_*^{\IndCoh}} \ar[d]^{\Psi} &&  \IndCoh^*(T)\ar[d]^{\Psi} \\
\QCoh(S) \ar[rr]^{f_*} && \QCoh(T)
}
\]

\noindent and such that $f_*^{\IndCoh}(\Coh(S)) \subset \QCoh(T)^+$.
(Indeed, this functor is necessarily the ind-extension of 
$\Coh(S) \xar{f_*} \QCoh(T)^+ \simeq \IndCoh^*(T)^+ \subset \IndCoh^*(T)$.

In general, the functor $f_*^{\IndCoh}$ can\footnote{Unfortunately, \cite{methods}, which
is currently being revised, erroneously claims otherwise.}be pathological, i.e., it may fail to map $\IndCoh^*(S)^+$ into $\IndCoh^*(T)^+$.
This issue complicates the theory, compared to the finite type setting. 

\subsubsection{}

Suppose $f:S \to T$ as above is eventually coconnective (i.e., of finite
$\Tor$-amplitude). In this case,
$f^*:\QCoh(T) \to \QCoh(S)$ maps $\Coh(T)$ to $\Coh(S)$. By ind-extension,
we obtain a functor $f^{*,\IndCoh}:\IndCoh^*(T) \to \IndCoh^*(S)$.
It is immediate to see that it is left adjoint to $f_*^{\IndCoh}$.

Moreover, $f^{*,\IndCoh}$ is manifestly right $t$-exact. Therefore, 
$f_*^{\IndCoh}$ \emph{is} left $t$-exact in this case. 

\subsubsection{}

Suppose $f:S \to T$ is proper (e.g., an almost finitely presented
closed embedding). 

Then $f_*:\Coh(S) \to \QCoh(T)^+$ maps into $\Coh(T)$. Therefore,
the functor $f_*^{\IndCoh}$ admits a continuous right adjoint, 
that we denote $f^!$ in this case. 

\subsubsection{Semi-coherence}

We now provide a general hypothesis that ensures that $f_*^{\IndCoh}$ is
left $t$-exact, even when $f$ is not eventually coconnective. 

\begin{defin}
\source{tates-thesis-de-rham-setting}
\notready

A connective and eventually coconnective 
commutative algebra $A$ is \emph{semi-coherent}
if every $M \in A\mod^+$ can be written as a filtered colimit
of objects $M_i$ such that:

\begin{itemize}

\item $M_i \in \Coh(A)$ % notation!!
for every $i$.

\item There exists an integer $r$ independent of $i$ 
such that every $M_i$ lies in $A\mod^{\geq -r}$.

\end{itemize}

We say $S = \Spec(A)$ is \emph{semi-coherent} if $A$ is.

\end{defin}

\begin{rem}
\source{tates-thesis-de-rham-setting}
\notready\label{r:semi-coh-ab}
\source{tates-thesis-de-rham-setting}

In the definition, if one restricts to $M \in A\mod^{\heart}$,
one finds that such an $r$ may be constructed (if it exists)
independent of $M$ (by considering direct sums of modules).  
It follows that it suffices to check the hypothesis just for $M \in A\mod^{\heart}$.

\end{rem}

\begin{example}
\source{tates-thesis-de-rham-setting}
\notready

Any Noetherian (and eventually coconnective) 
$A$ is obviously semi-coherent. More generally, any 
coherent (and eventually coconnective) $A$ is semi-coherent.

\end{example}

We now have the following simple observation.

\begin{lem}
\source{tates-thesis-de-rham-setting}
\notready\label{l:semi-coh-functor}
\source{tates-thesis-de-rham-setting}

Suppose $S$ is a semi-coherent eventually connective affine scheme.
Suppose $\sC \in \DGCat_{cont}$ is a DG category with a $t$-structure
compatible with filtered colimits.  
Suppose $F:\IndCoh^*(S) \to \sC$ is a continuous left $t$-exact DG functor 
such that $F(\Coh(S)) \subset \sC^+$.
Then $F(\IndCoh^*(S)^+) \subset \sC^+$.

\end{lem}

\begin{proof}

Suppose $\sF \in \IndCoh^*(S)^+$. By assumption and $S$, we can write
$\sF$ as a filtered colimit $\colim_i \sF_i$ with 
$\sF_i \in \Coh(S) \cap \QCoh(S)^{\geq -r}$ for
some $r$; equivalently, $\sF_i \in \Coh(S) \cap \IndCoh^*(S)^{\geq -r}$.
Then $F(\sF) = \colim_i F(\sF_i) \in \sC^{\geq -r}$ by assumption, giving
the claim.

\end{proof}

\begin{cor}
\source{tates-thesis-de-rham-setting}
\notready\label{c:semi-coh-push}
\source{tates-thesis-de-rham-setting}

Suppose $f:S \to T$ is a morphism of eventually coconnective affine schemes
with $S$ semi-coherent. Then $f_*^{\IndCoh}$ is $t$-exact.

\end{cor}

\begin{proof}

By Lemma \ref{l:semi-coh-functor}, $f_*^{\IndCoh}$ is left $t$-exact.
Therefore, we have a commutative diagram:
%
\[
\xymatrix{
\IndCoh^*(S)^+ \ar[rr]^{f_*^{\IndCoh}} \ar[d]^{\Psi} &&  \IndCoh^*(T)^+\ar[d]^{\Psi} \\
\QCoh(S)^+ \ar[rr]^{f_*} && \QCoh(T)^+
}
\]

\noindent in which the left arrows are $t$-exact equivalences and the
bottom arrow is $t$-exact (by affineness). This gives the claim.

\end{proof}

\subsubsection{}

We have the following important technical result.

\begin{prop}
\source{tates-thesis-de-rham-setting}
\notready\label{p:semi-coh}
\source{tates-thesis-de-rham-setting}

The (classical) commutative algebra $A_n$ defined in \S \ref{ss:coords} is
semi-coherent.

\end{prop}

In other words, $\widetilde{\sZ}^{\leq n}$ is semi-coherent.

We defer the proof of the proposition to \S \ref{ss:z-semi-coh-pf}, 
at the end of this section. Until that point, we assume it.

\subsection{Coherent sheaves with support}

\subsubsection{}

We now consider the following situation. 

Let $Z$ be an affine, eventually coconnective scheme. 
Let $f:Z \to \bA^1$ be a function. Let $U = \{f \neq 0\} \subset Z$, 
and let $j:U \to Z$ be the embedding.
Let $i:Z_0 \into Z$ be a closed subscheme with 
$Z \setminus Z_0 = U$.

The following type of result is standard in the finite type setting.

\begin{lem}
\source{tates-thesis-de-rham-setting}
\notready\label{l:coh-supp}
\source{tates-thesis-de-rham-setting}

Suppose that:

\begin{itemize}

\item $i$ is almost finitely presented.

\item The map $i_*^{\IndCoh}:\IndCoh^*(Z_0) \to \IndCoh^*(Z)$ is $t$-exact.

\end{itemize}

\noindent Then the functor:
%
\[
(j^{*,\IndCoh},i^!):\IndCoh^*(Z) \to \IndCoh^*(U) \times \IndCoh^*(Z_0)
\]

\noindent is conservative. 

\end{lem}

\begin{proof}

Let $\sF \in \IndCoh^*(Z)$ be given. We suppose $\sF$ is a non-zero
object with $j^{*,\IndCoh}(\sF) = 0$. Our goal is to construct 
an object $\sH \in \IndCoh^*(Z_0)$ and a non-zero map 
$i_*^{\IndCoh}(\sH) \to \sF$. 

\step We have:\footnote{Indeed, $j_*^{\IndCoh}$ is automatically left $t$-exact
(hence $t$-exact) because $j$ is flat; therefore, this formula follows from
the quasi-coherent setting by ind-extension.}
%
\[
j_*^{\IndCoh}j^{*,\IndCoh}(\sF) = \colim \big(\sF \xar{f} \sF \xar{f} \ldots \big) = 
0.
\]

Let $\sG \in \Coh(Z)$ be given with a non-zero map $\alpha:\sG \to \sF$; such 
$(\sG,\alpha)$
exists because $\sF$ is assumed non-zero. 

By compactness of $\sG$, there is an integer $n$ such that the composition:
%
\[
\sG \xar{\alpha} \sF \xar{f^n} \sF
\]

\noindent is null-homotopic. This map coincides with the composition:
%
\[
\sG \xar{f^n} \sG \xar{\alpha} \sF
\]

\noindent so this map must be null-homotopic as well. 
It follows that $\alpha$ factors through a (necessarily non-zero) map: 
%
\[
\sG/f^n\coloneqq \Coker(f^n:\sG \to \sG) \to \sF
\]

\noindent for some $n \gg 0$.

We have a standard co/fiber sequence:
%
\[
\sG/{f^{n-1}} \xar{f} \sG/f^n \to \sG/f.
\]

\noindent By descending induction, we see that there must exist
a non-zero map from $\sG/f \to \sF$. Replacing $\sG$ with $\sG/f \in \Coh(Z)$, 
we can assume that $f$ acts by zero on our original object $\sG$.

\step We now digress to make the following simple commutative algebra 
observation.\footnote{In our eventual application, $Z_0$ is the classical
scheme underlying $\{f=0\}$. This step is unnecessary under that
additional hypothesis.}
Suppose $Z = \Spec(A)$ below.

Let $I \subset H^0(A)$ be the ideal of 
elements vanishing on $Z_0^{cl}$. Let $(f) \subset H^0(A)$ 
be the ideal generated by $f$. We claim that there are integers $r,s$
such that:
%
\[
I^r \subset (f), (f^s) \subset I. 
\]

For any $g \in I$, we have $\{g \neq 0\} \subset \{f \neq 0\}$ by assumption. 
Therefore, $f$ is invertible in $H^0(A)[g^{-1}]$, so $g^{r_0} \in (f)$ for some
$r_0$ (depending on $g$). 
Because $i$ is almost finitely presented, $I$ is finitely generated,
so we obtain $I^r \subset (f)$ for $r \gg 0$.

The other inclusion (which we do not need) follows similarly: $\{f \neq 0\}$ is covered
by $\{g \neq 0\}$ for $g \in I$, so $1 \in I[f^{-1}] = H^0(A)[f^{-1}]$,
so $f^s \in I$ for some $s$.

We deduce the following: any module: 
%
\[
M \in H^0(A/f)\mod^{\heart} \subset H^0(A)\mod^{\heart} = 
A\mod^{\heart}
\]

\noindent admits a finite filtration:
%
\[
0 = I^r M \subset I^{r-1} M \subset \ldots \subset IM \subset M
\]

\noindent with subquotients lying in $H^0(A)/I\mod^{\heart} \subset A\mod^{\heart}$.

\step We now return to earlier setting.

Recall that we have $\sG \in \Coh(Z)$ on which $f$ acts null-homotopically.
Because $\sG$ is bounded, it has a finite Postnikov filtration as an object
of $\QCoh(Z)$
with associated graded terms $H^i(\sG)[-i]$. 
Each $H^i(\sG) \in 
\QCoh(Z)^{\heart}$. Because $f$ acts by zero on them, they lie
in $\QCoh(\{f = 0\})^{\heart}$. By the previous step, each
of these terms is obtained by successively extending 
objects in $\QCoh(Z_0)^{\heart}$.

Now the diagram:
%
\[
\xymatrix{
\IndCoh(Z_0)^+ \ar[r]^{i_*^{\IndCoh}} \ar[d]^{\Psi} & \IndCoh(Z)^+ \ar[d]^{\Psi} \\
\QCoh(Z_0)^+ \ar[r]^{i_*} & \QCoh(Z)^+}
\]

\noindent commutes because we assumed $i_*^{\IndCoh}$ $t$-exact;
moreover, the vertical arrows are equivalences. Therefore, $\sG$
admits a finite filtration as an object of $\IndCoh^*(Z)$
with associated graded terms $\IndCoh$-pushed forward from $Z_0$. 
Therefore, one of these associated graded terms must have a non-zero map
to $\sF$, giving the claim.  

\end{proof}

\subsubsection{Application to $\widetilde{\sZ}^{\leq n}$}

We now apply the above to construct generators of $\IndCoh^*(\widetilde{\sZ}^{\leq n})$.

\begin{prop}
\source{tates-thesis-de-rham-setting}
\notready\label{p:tilde-z-gens}
\source{tates-thesis-de-rham-setting}

For
$n > 0$, the DG category $\IndCoh^*(\widetilde{\sZ}^{\leq n})$ is compactly 
generated by the objects 
$\{\iota_*^{r,\IndCoh}(\sO_{\widetilde{\sZ}^{\leq n}})\}_{r \geq 0}$.

For $n = 0$, the DG category $\IndCoh^*(\widetilde{\sZ}^{\leq 0})$ is compactly 
generated by $\sO_{\widetilde{\sZ}^{\leq 0}})$.

\end{prop}

\begin{proof}

We proceed by induction on $n$. 

\step 

The $n = 0$ case is simple: $\widetilde{\sZ}^{\leq 0}$ is $\bA^1$ (with coordinate
$a_0$). 

\step 

We now fix $n > 0$, and that the proposition for $n - 1$. 
We introduce the following notation.   

Let $\sC \subset \IndCoh^*(\widetilde{\sZ}^{\leq n})$ be the subcategory generated
by the objects $\iota_*^{r,\IndCoh}\sO_{\widetilde{\sZ}^{\leq n}}$. We need
to show that $\sC$ is the full $\IndCoh^*$.

\step \label{st:gen-3}
\source{tates-thesis-de-rham-setting}

We observe that \eqref{eq:fund-ses} implies that 
$\delta_{n,*}^{\IndCoh}(\sO_{\sZ}^{\leq n-1}) \in \sC$. 
More generally, we find 
$\iota_*^{r,\IndCoh}\delta_{n,*}^{\IndCoh}(\sO_{\sZ}^{\leq n-1}) \in \sC$ for
any $r\geq 0$. 

\step 

Below, we will apply Lemma \ref{l:coh-supp} with 
$i:Z_0 \to Z$ corresponding to 
$\delta_n:\widetilde{\sZ}^{\leq n-1} \to \widetilde{\sZ}$. We remark
that $\delta_n$ is almost finitely presented by 
Lemma \ref{p:z-flat-axes} \eqref{i:z-axes-2} 
(cf. the proof of Corollary \ref{c:iota-afp}), and that $\delta_{n,*}^{\IndCoh}$
is $t$-exact by Proposition \ref{p:semi-coh}. 

In other words, our application of the lemma is justified. 

\step 

We now treat the $n = 1$ case. 

Suppose $\sF \in \IndCoh^*(\widetilde{\sZ}^{\leq 1})$ lies in the 
right orthogonal to $\sC$. 

By Step \ref{st:gen-3} and the $n = 0$ case, $\delta_1^!(\sF) = 0$, and 
$\delta_1^!(\iota^r)^! \sF = 0$ more generally. Therefore, 
by Lemma \ref{l:coh-supp}, $\sF$ is $*$-extended from:
%
\[
\widetilde{\sZ}^{\leq 1} \setminus  
\big(\widetilde{\sZ}^{\leq 0} \cup \ldots \cup \iota^r(\widetilde{\sZ}^{\leq 0})\big)
\]

\noindent for each $r$. It follows that $\sF$ is $*$-extended from:
%
\[
\widetilde{\sU}_1 \coloneqq 
\widetilde{\sZ}^{\leq 1} \times_{\sZ^{\leq 1}} \sU_1.
\]

But by Lemma \ref{l:u1} (cf. \S \ref{sss:u-reg-noeth}), 
$\IndCoh^*(\widetilde{\sU}_1) \simeq \QCoh(\widetilde{\sU}_1)$ 
is generated by its structure sheaf. Since this is the $*$-restriction
of the structure sheaf on $\sZ^{\leq 1}$, we obtain the claim. 

\step 

We now assume $n>1$ case. 

In this case, we have:
%
\[
\iota_*^{r,\IndCoh}\delta_{n,*}^{\IndCoh} \simeq 
\delta_{n,*}^{\IndCoh}
\iota_*^{r,\IndCoh}.
\]

\noindent (Here we use $\iota$ to denote the map for both $\sZ^{\leq n}$ 
and $\sZ^{\leq n-1}$.)

By induction and Step \ref{st:gen-3}, $\sC$ contains the subcategory
generated under colimits by the essential image of
$\delta_{n,*}^{\IndCoh}$. Therefore, by Lemma \ref{l:coh-supp}, 
any $\sF$ in the right orthogonal to $\sC$ is $*$-extended
from $\widetilde{\sU}_n$. 

But again, $\IndCoh^*(\widetilde{\sU}_n) \simeq \QCoh(\widetilde{\sU}_n)$ is generated
by its structure sheaf by Lemma \ref{l:un}; as 
$\sO_{\widetilde{\sZ}^{\leq n}} \in \sC$, we obtain the claim.

\end{proof}


\subsection{Stacky case}

\subsubsection{}

Now suppose $S$ is a stack of the form $T/G$ for $G$ classical affine group scheme
and $T$ an eventually coconnective affine scheme.

In this case, $G$ acts weakly on $\IndCoh^*(T)$, i.e., $\QCoh(G)$
acts canonically on $\IndCoh^*(T)$ (by flatness of $G$ over a point). 
We set $\IndCoh^*(S) \coloneqq \IndCoh^*(T)^{G,w}$.
%One can see that this is canonically independent of the presentation of
%$S$ as $T/G$, cf. \cite{methods} Corollary 6.25.2.

%We say an object $\sF \in \IndCoh^*(S)$ is \emph{coherent} if
%its pullback to some (equivalently any) cover is coherent. 
%We let $\Coh(S) \subset \IndCoh^*(S)$ denote the corresponding
%full subcategory. We remark that in general, objects of
%$\Coh(S)$ may not be compact, but this is not an issue if
%$G$ is assumed to be of finite type. 

There is a unique $t$-structure on $\IndCoh^*(S)$ such that the
forgetful functor:
%
\[
\IndCoh^*(S) \to \IndCoh^*(T)
\]

\noindent is $t$-exact. 

\subsubsection{}

In the above setting, given a map $f:S_1 \to S_2$ of stacks of the
above types, we obtain similar functoriality as in the non-stacky case;
we omit the details.  

\subsubsection{}

We now have:

\begin{cor}
\source{tates-thesis-de-rham-setting}
\notready\label{c:z-gen}
\source{tates-thesis-de-rham-setting}

The category $\IndCoh^*(\sZ^{\leq n})$ is compactly generated
by the objects $\iota_*^{r,\IndCoh}(\sO_{\sZ^{\leq n}})(m)$, where $r \geq 0$ and 
$m \in \bZ$.

\end{cor}

\begin{proof}

For $G$ an affine algebraic group (in particular, of finite type) and 
$\sF \in \sC^{G,w}$, it is a general fact that $\sF$ is compact
if and only if its image $\Oblv(\sF) \in \sC$ is compact.\footnote{See 
\cite{methods}, proof of Lemma 5.20.2.}
Moreover, the functor $\Oblv:\sC^{G,w} \to \sC$ always generates
the target under colimits. 

Therefore, writing $\sZ^{\leq n} = \widetilde{\sZ}^{\leq n}/\bG_m$, 
the above objects are indeed compact. The generation follows from 
Proposition \ref{p:tilde-z-gens}.

\end{proof}

\subsubsection{}

There is a natural functor:
%
\[
\Psi:\IndCoh^*(\sZ^{\leq n}) \to \QCoh(\sZ^{\leq n})
\]

\noindent obtained by passing to weak $\bG_m$-invariants for
the corresponding functor for $\widetilde{\sZ}^{\leq n}$. This functor
is an equivalence on coconnective (equivalently: eventually coconnective) 
subcategories via the corresponding assertion for $\widetilde{\sZ}^{\leq n}$.

We let $\Coh(\sZ^{\leq n}) \subset \IndCoh^*(\sZ^{\leq n})$ 
denote the subcategory of compact objects. By the above, 
the functor $\Coh(\sZ^{\leq n}) \to \QCoh(\sZ^{\leq n})$
induced by $\Psi$ is fully faithful; its essential image is 
exactly the subcategory of 
eventually coconnective almost compact objects. 

\subsubsection{}

We now observe that the functor:
%
\[
\iota_*^{\IndCoh}:\IndCoh^*(\sZ^{\leq n}) \to \IndCoh^*(\sZ^{\leq n})
\]

\noindent is $t$-exact. 

Indeed, by definition, this assertion reduces to the analogous
statement for $\widetilde{\sZ}^{\leq n}$, which in turn follows 
from Proposition \ref{p:semi-coh}.

The same analysis applies for pushforward functors:
%
\[
\delta_{n+1,*}^{\IndCoh}:\IndCoh^*(\sZ^{\leq n}) \to \IndCoh^*(\sZ^{\leq n+1}).
\]

\subsection{Ind-algebraic stacks}

We define:
%
\[
\IndCoh^*(\sY_{\log}^{\leq n}) \coloneqq \colim \, \big(
\IndCoh^*(\sZ^{\leq n}) \xar{\iota_*^{\IndCoh}}
\IndCoh^*(\sZ^{\leq n}) \xar{\iota_*^{\IndCoh}} \ldots \big) \in \DGCat_{cont}.
\]

\noindent We let: 
%
\[
i_{n,*}^{\IndCoh}: \IndCoh^*(\sZ^{\leq n}) \to \IndCoh^*(\sY_{\log}^{\leq n})
\]

\noindent denote the structural functor. 

Similarly, we define:
%
\[
\IndCoh^*(\sY_{\log}) \coloneqq \underset{n}{\colim} \, \IndCoh^*(\sY_{\log}^{\leq n}) =
\underset{\iota_*^{\IndCoh},\delta_{n,*}^{\IndCoh}}{\underset{n,m}{\colim}} \, \IndCoh^*(\sZ^{\leq n}) \in \DGCat_{cont}.
\]

By the above and \cite{whit} Lemma 5.4.3, there is a unique 
$t$-structure on $\IndCoh^*(\sY_{\log}^{\leq n})$ such that 
each structural functor:
%
\[
\IndCoh^*(\sZ^{\leq n}) \to \IndCoh^*(\sY_{\log}^{\leq n})
\]

\noindent is $t$-exact. Similarly, there is a unique $t$-structure on 
$\IndCoh^*(\sY_{\log})$ such that each structural functor:
%
\[
\IndCoh^*(\sY_{\log}^{\leq n}) \to \IndCoh^*(\sY_{\log})
\]

\noindent is $t$-exact.

Each of these categories is compactly generated as all of our structural
functors preserve compact objects.

\subsection{Grassmannian actions}

Suppose $n>0$ in what follows. 

\subsection{}

Recall that $\Gr_{\bG_m}$ tautologically acts on $\sY_{\log}$.

Observe that $\sZ \subset \sY_{\log}$ is preserved
under the action of the submonoid 
$\Gr_{\bG_m}^{pos} \subset \Gr_{\bG_m}$.

Indeed, suppose more generally that $Y$ is a prestack
mapping to $\fL \bA^1/\bG_m$. Then:
%
\[
\Gr_{\bG_m}^{pos} = \Ker(\fL^+ \bA^1/\bG_m \to \fL \bA^1/\bG_m)
\]

\noindent acts canonically on:
%
\[
Z \coloneqq Y \underset{\fL \bA^1/\bG_m}{\times} \fL^+ \bA^1/\bG_m
\]

\noindent compatibly with the canonical $\Gr_{\bG_m} = 
\Ker(\fL^+\bB \bG_m \to \fL \bB \bG_m)$-action on:
%
\[
Y_{\log} \coloneqq Y \underset{\fL\bB\bG_m}{\times} \fL^+ \bB \bG_m.
\]

Taking $Y = \sY$, we recover the claim.

\subsubsection{}

Next, observe that the action of 
$\Gr_{\bG_m}^{pos,\leq n} \subset \Gr_{\bG_m}^{\leq n}$
on $\sY_{\log}^{\leq n}$ preserves $\sZ^{\leq n}$.

Indeed, this is an immediate consequence of the above.

\subsubsection{}

By the above, there is a canonical action:
%
\[
\QCoh(\Gr_{\bG_m}^{pos,\leq n}) \actson \QCoh(\sZ^{\leq n})
\]

\noindent where the left hand side is given its convolution monoidal
structure. 

For $\sF \in \Coh(\Gr_{\bG_m}^{pos,\leq n}) \subset \QCoh(\Gr_{\bG_m}^{\leq n})$,
the corresponding action functor:
%
\[
\QCoh(\sZ^{\leq n}) \to \QCoh(\sZ^{\leq n})
\]

\noindent clearly preserves $\Coh(\sZ^{\leq n})$. Therefore, we obtain 
a canonical action:
%
\[
\IndCoh(\Gr_{\bG_m}^{pos,\leq n}) \actson \IndCoh^*(\sZ^{\leq n}).
\]

\subsubsection{}

By Remark \ref{r:gr-n-pos-gp-cpl}, we have:
%
\[
\IndCoh(\Gr_{\bG_m}^{\leq n}) \underset{\IndCoh(\Gr_{\bG_m}^{pos,\leq n})}{\otimes} 
\simeq \big(
\IndCoh^*(\sZ^{\leq n}) \xar{\iota_*^{\IndCoh}}
\IndCoh^*(\sZ^{\leq n}) \xar{\iota_*^{\IndCoh}} \ldots \big) \simeq 
\IndCoh^*(\sY_{\log}^{\leq n}).
\]

\noindent Thus, we obtain a natural action of $\IndCoh(\Gr_{\bG_m}^{\leq n})$
on $\IndCoh^*(\sY_{\log}^{\leq n})$ such that the functor
$i_{n,*}^{\IndCoh}$ is a morphism of $\IndCoh(\Gr_{\bG_m}^{pos,\leq n})$-module
categories. 

\subsubsection{}

The map $\delta_{n,*}^{\IndCoh}:\IndCoh^*(\sZ^{\leq n}) \to \IndCoh^*(\sZ^{\leq n+1})$
is naturally a map of $\IndCoh(\Gr_{\bG_m}^{pos,\leq n})$-module categories,
as is evident by again considering subcategories of compact objects. 

Therefore, $\IndCoh(\Gr_{\bG_m})$ naturally acts on $\IndCoh(\sY_{\log})$,
compatibly with the above actions. 

\subsection{Ind-coherent sheaves on $\sY$}

We now define:
%
\[
\IndCoh^*(\sY) \coloneqq \IndCoh^*(\sY_{\log})^{\Gr_{\bG_m},w}.
\]

\noindent By definition, the right hand side is:
%
\[
\TwoHom_{\IndCoh(\Gr_{\bG_m})\mod}(\Vect, \IndCoh^*(\sY_{\log})).
\]

We let:
%
\[
\Oblv:\IndCoh^*(\sY) \to \IndCoh^*(\sY_{\log})
\]

\noindent denote the (conservative) forgetful functor.
It admits a left adjoint:
%
\[
\Av_!^{\Gr_{\bG_m},w}.
\] 

The resulting monad on $\IndCoh^*(\sY_{\log})$ is easily seen to be 
$t$-exact.\footnote{This reduces to the assertion that for 
$\sF \in \IndCoh(\Gr_{\bG_m}^{\leq n,pos})^{\heart}$, the 
action functor $\IndCoh^*(\sZ^{\leq n}) \to \IndCoh^*(\sZ^{\leq n})$
is $t$-exact. Filtering $\sF$, we are reduced to the case that
it is the skyscraper sheaf at a $k$-point. In that case, the
relevant functor is a composition of functors $\iota_*^{\IndCoh}$,
so $t$-exact by Proposition \ref{p:semi-coh}.}
Therefore, we obtain:

\begin{prop}
\source{tates-thesis-de-rham-setting}
\notready\label{p:y-t-str}
\source{tates-thesis-de-rham-setting}

There is a unique $t$-structure on $\IndCoh^*(\sY)$ such that the
functor $\Oblv:\IndCoh^*(\sY) \to \IndCoh^*(\sY_{\log})$ is $t$-exact.
Moreover, the functor $\Av_!^{\Gr_{\bG_m},w}$ is $t$-exact.

\end{prop}

Indeed, this follows from:

\begin{lem}
\source{tates-thesis-de-rham-setting}
\notready

Let $\sC \in \DGCat_{cont}$ be equipped with a $t$-structure compatible with
filtered colimits and a right $t$-exact monad $T:\sC \to \sC$.
Then $T\mod(\sC)$ admits a unique $t$-structure such that the
forgetful functor $T\mod(\sC) \to \sC$ is $t$-exact. 

\end{lem}

\subsubsection{}\label{ss:y-trun-indcoh}
\source{tates-thesis-de-rham-setting}

In the truncated setting, we similarly define:
%
\[
\IndCoh^*(\sY^{\leq n}) \coloneqq 
\IndCoh^*(\sY_{\log}^{\leq n})^{\Gr_{\bG_m}^{\leq n},w}.
\]

\noindent We again have adjoint functors:
%
\[
\Oblv,\Av_!^{\Gr_{\bG_m}^{\leq n},w}
\]

\noindent and a $t$-structure on $\IndCoh^*(\sY^{\leq n})$.

\subsubsection{}

By functoriality of the constructions, there is a natural functor:
%
\[
\lambda_{n,*}^{\IndCoh}:\IndCoh^*(\sY^{\leq n}) \to \IndCoh^*(\sY).
\]

\noindent This functor is $t$-exact and preserves compact objects. Moreover, we have:

\begin{lem}
\source{tates-thesis-de-rham-setting}
\notready\label{l:lambda-ff}
\source{tates-thesis-de-rham-setting}

For $n>0$, the functor $\lambda_{n,*}^{\IndCoh}$ is fully faithful.

\end{lem}

Roughly speaking, this is true because the map $\sY^{\leq n} \to \sY$ is 
formally \'etale for $n>0$, which follows from the corresponding
fact for $\LocSys_{\bG_m}^{\leq n} \to \LocSys_{\bG_m}$ (which 
follows from Proposition \ref{p:locsys-log}). Unwinding the constructions
to convert this argument into a proof is straightforward. 

\subsection{Proof of Proposition \ref{p:semi-coh}}\label{ss:z-semi-coh-pf}
\source{tates-thesis-de-rham-setting}

We now give the proof of Proposition \ref{p:semi-coh}, which we deferred earlier.

\subsubsection{}

Our argument is by induction on $n$. We proceed in steps.

\subsubsection{Step 1: $n = 0$ case}

The $n = 0$ case is trivial, as $A_0 = k[a_0]$ is Noetherian.

\subsubsection{Step 2: $n = 1$ case} 

The inductive step we give for $n>1$ below may be adapted to treat the 
$n = 1$ case, but it requires somewhat more work; we indicate how this works
in \S \ref{sss:semi-coh-7}. We prefer
to give a direct argument. Actually, we will show that $A_1$ is coherent.

Recall the Noetherian subalgebras $A_{1,\leq m} \subset A_1$ from \S \ref{ss:rs-sketch}.
It suffices to show that for any finitely presented $M \in A_1\mod^{\heart}$,
there is an integer $m \gg 0$ and finitely generated
$N \in A_{1,\leq m}\mod^{\heart}$ such that
there is an isomorphism $N \otimes_{A_{1,\leq m}} A_1 \isom M$, where
the tensor product is understood in the derived sense.

Choose some $m_0$ and $N_0 \in A_{1,\leq m_0}\mod^{\heart}$ with an 
isomorphism $H^0(N_0 \otimes_{A_{1,\leq m_0}} A_1) \isom M$; this may be
done by choosing generators and relations for $M$ and $m_0$ such that 
the relations all involve linear combinations with coefficients in $A_{1,\leq m_0}$.

Choose $m > m_0$ such that $N_0$ does not contain any $(b_{-1}-i)$-torsion 
for any integer $i\geq m$; this can be done because the finitely generated
$A_{1,\leq m_0}$-module $N_0$ only admits finitely many associated primes.
Then it is easy to see that $N \coloneqq 
H^0(A_{1,\leq m} \otimes_{A_{1,\leq m_0}} N_0)$ satisfies the hypotheses.

\begin{rem}
\source{tates-thesis-de-rham-setting}
\notready

We see here that $A_1$ is actually coherent, not merely semi-coherent.
We do not consider this question for general $n$.

\end{rem}

\subsubsection{Step 3: filtered colimits, extensions, and effective bounds} 
We now make the following observations for d\'evissage.

For an integer $r$, we say that $M \in A_n\mod$ is \emph{$r$-good} if $M$ can be 
expressed as a filtered colimit of coherent objects concentrated in 
degrees $\geq -r$. Clearly any such $M$ lies in $A_n\mod^{\geq -r}$.

In other words, $M$ is $r$-good if it lies in the full subcategory:
%
\[
\Ind(\Coh(A) \cap A\mod^{\geq -r}) \subset A\mod^{\geq -r}.
\]

\noindent Here the natural functor is fully faithful because
$\Coh(A) \cap A\mod^{\geq -r}$ consists of objects that are 
compact in $A\mod^{\geq -r}$ by assumption.

It follows from the second description that $r$-good objects are closed
under filtered colimits (for fixed $r$). We also observe that $r$-good objects
are closed under extensions.

Finally, we say an object $M \in A_n\mod$ is \emph{good} if it is
$r$-good for some $r$. So our task is to show that any $M \in A_n\mod^+$ 
is good, or equivalently, any $M \in A_n\mod^{\heart}$.

\subsubsection{Step 4: reductions}

We now begin our induction. Take $n>1$ and assume the result for $n-1$.

Suppose $M \in A_n\mod^+$. 
It it suffices to show that:
%
\[
M[b_{-n}^{-1}] \coloneqq \underset{b_{-n}}{\colim} \, M \text{ and }
\Coker(M \to M[b_{-n}^{-1}])
\]

\noindent are good (where $\Coker$ indicates the \emph{homotopy cokernel}, 
i.e., the cone).

We check these assertions below. We remark that by Remark \ref{r:semi-coh-ab},
we are reduced to considering $M \in A_n\mod^{\heart}$.

\subsubsection{Step 5: generic case}

Note that $M[b_{-n}^{-1}] \in A_n[b_{-n}^{-1}]\mod^{\heart} \subset A_n\mod^{\heart}$.

As $n>1$, we have an isomorphism:
%
\[
A_n[b_{-n}^{-1}] = k[b_{-1},\ldots,b_{-n},b_{-n}^{-1}].
\]

\noindent by Lemma \ref{l:un}.

As this algebra is regular Noetherian, it follows that
$M[b_{-n}^{-1}]$ has bounded $\Tor$-amplitude as an 
$A_n[b_{-n}^{-1}]$-module (simply because it is bounded from below),
so the same is true of $M[b_{-n}^{-1}]$ as an $A_n$-module.
By standard homological algebra, $M[b_{-n}^{-1}]$ can be represented
by a bounded complex of flat (classical) $A_n$-modules. 
By Lazard, a flat $A_n$-module is good, so $M[b_{-n}^{-1}]$ is good.

\subsubsection{Step 6: torsion case} 

We now show that $\widetilde{M} = \Coker(M \to M[b_{-n}^{-1}])$ is good. 
More generally, we show that any bounded complex 
$\widetilde{M} \in A_n\mod^+$ such that $b_{-n}$ acts locally nilpotently on
its cohomologies is good. As the complex $\widetilde{M}$ is assumed bounded,
it suffices to treat its cohomology groups one at a time, so 
we can assume (up to shifting) that $\widetilde{M} \in A_n\mod^{\heart}$.

By induction, there exists an integer $r$ such that any module
$N \in A_{n-1}\mod^{\heart}$ is $r$-good as an $A_{n-1}$-module.
Because the projection $A_n \xar{b_{-n} \mapsto 0} A_{n-1}$ is
almost finitely presented by Proposition \ref{p:z-flat-axes} \eqref{i:z-axes-2},
so coherent $A_{n-1}$-complexes restrict to coherent $A_n$-complexes, 
$N$ restricts to an $r$-good $A_n$-module. 

Now let $\widetilde{M}_i \subset \widetilde{M}$ be the 
non-derived kernel of the map $b_{-n}^i:\widetilde{M} \to \widetilde{M}$. 
Clearly $\widetilde{M}_i/\widetilde{M}_{i-1} \in A_{n-1}\mod^{\heart} \subset
A_n\mod^{\heart}$, so is $r$-good for the above $r$ (which is independent
of everything in sight except $n$). The module 
$\widetilde{M}_i$ is then $r$-good, as it is obtained by successively extending
$r$-good modules. Finally, $\widetilde{M} = \colim_i \widetilde{M}_i$ is 
a filtered colimit of $r$-good modules, so is $r$-good. This concludes the argument.

\subsubsection{Step 7: revisiting the $n = 1$ case}\label{sss:semi-coh-7}
\source{tates-thesis-de-rham-setting}

Finally, we make a remark that the above argument can be adapted to 
treat the $n = 1$ case, if one so desires. As in Lemma \ref{l:u1}, 
we should consider the localization of 
$A_1$ at $\{b_{-1},b_{-1}-1,b_{-1}-2,\ldots\}$
(instead of simply at $b_{-1}$); one again then obtains a Noetherian ring
of finite global dimension. The argument concludes by noting that 
any classical $A_1/(b_{-1}-i) = k[a_i]$-module is $0$-good, so 
any module for $A_1$ that is torsion for the above elements (thus
necessarily the direct sum of its torsion with respect to each) is
$0$-good (by an argument as above).

In other words, we need to localize at infinitely many elements before
obtaining a regular ring; the saving
grace is that the goodness for torsion modules at each is bounded independently
of the element.
